\subsection{Propagation Mechanics}

When Python code raises an exception, CPython does not rely on hardware traps or portable \texttt{setjmp}/\texttt{longjmp} scaffolding in user code. The interpreter loop consults a static \emph{exception table} baked into each code object at compile time and rewires the virtual evaluation stack accordingly.

\subsubsection{Exception Routing, Stack Unwinding, and the Runtime Exception Table}

Modern CPython (3.11+) stores handler metadata in \texttt{co\_exceptiontable}, a compact byte sequence parallel to \texttt{co\_code}. Each entry maps a bytecode offset range to:

\begin{itemize}
    \item the handler target offset (where execution resumes after a match);
    \item the stack depth to restore;
    \item the exception class filter (if any) for \texttt{except} clauses;
    \item flags for \texttt{finally}, \texttt{with}, and loop-\texttt{else} cleanup regions.
\end{itemize}

When an exception is raised, the eval loop:

\begin{enumerate}
    \item captures the active exception state in thread-local storage;
    \item walks the current frame's instruction pointer backward through the exception table;
    \item finds the innermost entry covering the active offset;
    \item unwinds the value stack to the recorded depth;
    \item jumps to the handler label---\texttt{except} body, \texttt{finally} prologue, or \texttt{WITH} exit dispatch.
\end{enumerate}

If no handler exists in the current frame, the frame is popped, a traceback link is appended, and the search repeats in the caller's \texttt{co\_exceptiontable}. This continues until a match is found or the interpreter boundary prints an uncaught traceback.

This is a \emph{zero-cost} exception model in the steady state: no per-\texttt{try} frame allocation at runtime. The cost is paid at compile time when the table is constructed; the hot path is a table lookup plus stack pointer adjustment, not OS signal delivery.

Contrast with older CPython (pre-3.11) block-stack metadata and with C error codes:

\begin{verbatim}
C error-code model:  return -1; caller must poll every boundary
CPython EAFP model:  raise -> table lookup -> jump to handler label
\end{verbatim}

Nested \texttt{try}/\texttt{finally}/\texttt{with} structures compile to overlapping table entries ordered by specificity. An inner \texttt{finally} must run before an outer handler observes a propagated exception---the table encodes that ordering as sequential unwind targets, not as runtime stack walks searching for destructors.

Explicit \texttt{raise} re-enters the same mechanism: it sets the exception state and triggers an immediate table lookup from the current instruction offset, creating a deliberate non-local edge in the control-flow graph.

Traceback objects record the chain of code objects and line numbers visited during unwinding; they are diagnostic artifacts attached to the exception, not the routing mechanism itself. The router is \texttt{co\_exceptiontable}; the traceback is the audit log of frames traversed when no local handler matched.

Re-raising with bare \texttt{raise} preserves the active exception context while continuing propagation outward. \texttt{raise ... from ...} links \texttt{\_\_cause\_\_} for chained diagnostics but does not change the unwind algorithm: the exception table still drives frame exit and handler selection.

\begin{quote}
Exceptions in CPython are compile-time annotated control-flow edges. \texttt{co\_exceptiontable} is the jump ledger the interpreter consults to unwind the virtual stack and reroute execution without hardware fault semantics.
\end{quote}
